% kernel/part_intro-DE.tex

\chapter*{Ingenieurbegriffe}
\phantomsection
\addcontentsline{toc}{chapter}{Ingenieurbegriffe}

Die Berlinish-Grammatik hat eine Familie von Übergangskonstruktionen ausdrückbar
gemacht. Angenommen, axtranNew oder ein anderer Berechnungsdienst wählt nun eine
Eingangshalbwelle, einen Klothoidenkern und eine Ausgangshalbwelle und liefert
Parameter, welche die angegebenen geometrischen Randbedingungen erfüllen.
Innerhalb dieser Berechnung ist das Ergebnis mathematisch gültig. Es ist damit
noch nicht der angenommene Übergang eines Eisenbahn-Alignments.

Damit das Ergebnis als Ingenieurinformation nutzbar wird, muss es eine Bedeutung
erhalten, die weder Formel noch Solver bereitstellen können. Es muss sich auf ein
identifizierbares Alignment und einen bestimmten Zustand dieses Alignments
beziehen. Seine geordnete Konstruktion muss prüfbar bleiben. Metrische
Realisierung, Koordinatenreferenzsystem, Interpretation der Stationierung,
Quelle, Ableitung und zeitliche Qualifikation müssen bekannt sein. Danach ist
das Ergebnis für einen angegebenen Zweck zu bewerten, etwa für den Betrieb mit
einer bestimmten Geschwindigkeit unter anwendbaren Überhöhungsbedingungen.
Schließlich muss eine befugte Autorität---nicht die Berechnung---entscheiden, ob
es Ingenieurarbeit bestimmen darf.

Abbildung~\ref{fig:transition-meaning-chain} zeigt den Weg durch diesen Teil. Die
Kapitel führen keine Reihe unabhängiger Container ein. Sie verfolgen einen
Übergangsvorschlag, während seine mathematische Möglichkeit zu einer Aussage
über ein Engineering Object und danach zu Evidenz für einen anwendbaren und
möglicherweise angenommenen Zustand wird.

\begin{figure}[htbp]
  \centering
  \resizebox{\textwidth}{!}{% figures/kernel/transition_meaning_chain-DE.tex
\begin{tikzpicture}[
  >=Stealth,
  node distance=0.38cm,
  stage/.style={draw=black!55, rounded corners=2pt, fill=black!2,
    minimum width=2.15cm, minimum height=1.35cm, align=center,
    font=\scriptsize, inner sep=4pt},
  flow/.style={->, semithick, draw=black!65},
  status/.style={font=\scriptsize\itshape, text=black!72, align=center}
]
  \node[stage] (grammar) {Berlinish-\\Komposition\\$H_{in}+C+H_{out}$};
  \node[stage, right=of grammar] (calc) {parametrisierter\\Berechnungs-\\vorschlag};
  \node[stage, right=of calc] (object) {Alignment-Objekt,\\Zustand und\\Konstruktion};
  \node[stage, right=of object] (real) {metrische\\Realisierung und\\Koordinaten};
  \node[stage, right=of real] (evidence) {Repräsentation,\\Beobachtung,\\Provenienz};
  \node[stage, right=of evidence] (eval) {zweckrelative\\Bewertung und\\Anwendbarkeit};
  \node[stage, right=of eval] (decision) {autorisierte\\Ingenieur-\\entscheidung};
  \draw[flow] (grammar) -- (calc);
  \draw[flow] (calc) -- (object);
  \draw[flow] (object) -- (real);
  \draw[flow] (real) -- (evidence);
  \draw[flow] (evidence) -- (eval);
  \draw[flow] (eval) -- (decision);
  \node[status, below=0.18cm of calc] {berechnet};
  \node[status, below=0.18cm of eval] {für einen Zweck\\anwendbar};
  \node[status, below=0.18cm of decision] {innerhalb bestimmter\\Autorität angenommen};
\end{tikzpicture}
}
  \caption{Ein Übergangsvorschlag erhält durch ausdrückliche Identität,
  Konstruktion, Realisierung, Provenienz, Bewertung und Autorität
  Ingenieurbedeutung. Kein Pfeil verleiht stillschweigend den Status der
  nächsten Stufe.}
  \label{fig:transition-meaning-chain}
\end{figure}

Die leitende Unterscheidung lautet daher
\begin{equation}
\text{berechnet}\neq\text{angenommen}\neq\text{anwendbar}
\neq\text{autoritativ}.
\label{eq:calculation-status-chain}
\end{equation}
Ein berechneter Vorschlag kann die Bedingungen für Geschwindigkeit oder
Überhöhung verfehlen, von einer ungeeigneten Realisierung abhängen, auf
unvollständiger Evidenz beruhen oder nur eine Alternative unter mehreren
bleiben. Umgekehrt kann ein historisch angenommener Übergang autoritative
Evidenz für eine bestehende Anlage sein, ohne mathematisch optimal zu sein. Die
Ingenieurbegriffe bewahren diese Unterschiede und halten zugleich das
zugrunde liegende Alignment durch Änderungen hindurch erkennbar.
